\section{Internal Bell constraints from actual regional contractions}%
\label{sec:peps_vertex_bond_region_constraint}

The virtual-pair argument of
\cite[Section~IV.C.1]{Cirac2021Matrix}, local source lines 2017--2044,
uses the Bell condition on each original graph edge. Here this condition
is derived from the actual regional virtual bond maps, before any
physical injectivity assumption is imposed.

\begin{theorem}[Internal endpoint consistency and label symmetry]%
    \label{thm:peps_region_virtual_internal_bond}
    \lean{TNLean.PEPS.edgeLabelPerm,
        TNLean.PEPS.regionVertexEdgeLabelPerm,
        TNLean.PEPS.regionVirtualBondWeight_edgeLabelPerm,
        TNLean.PEPS.regionVirtualBondWeight_eq_zero_of_internal_off_diagonal,
        TNLean.PEPS.regionVirtualBondMap_range_edgeLabelPerm,
        TNLean.PEPS.regionVirtualBondMap_range_eq_zero_of_internal_off_diagonal}
    \leanok
    \uses{def:peps_injective_vertex_coordinates}
    Let both endpoints of an edge $e$ belong to a finite region $R$.
    Every vector in the actual regional virtual bond range vanishes
    when the two half-edge labels of $e$ differ. Its coefficients
    are unchanged when the same permutation of the bond alphabet
    is applied to both labels. The assertion includes arbitrary
    boundary coefficients and vanishing bond dimensions.
\end{theorem}

\begin{proof}\leanok
    Each incident-edge assignment uses one common label at the two
    endpoints of $e$, so no assignment contributes to a coefficient
    with unequal labels. A permutation of this common internal label
    is a bijection of the incident assignments, leaves every boundary
    label fixed, and permutes the two endpoint labels together.
    The resulting identities extend by linearity to the actual
    regional virtual map.
\end{proof}

\begin{theorem}[An internal edge satisfies every Bell slice condition]%
    \label{thm:peps_region_virtual_bell_slice}
    \lean{TNLean.PEPS.vertexVirtual_edgeLabelPerm_of_regionSlices,
        TNLean.PEPS.vertexVirtual_eq_zero_of_regionSlices_off_diagonal,
        TNLean.PEPS.fullRegionVertexConfigEquivEdgePair_edgeLabelPerm,
        TNLean.PEPS.edgePair_eq_zero_of_regionSlices_off_diagonal,
        TNLean.PEPS.edgePair_diagonal_update_eq_of_regionSlices,
        TNLean.PEPS.virtualBondSlice_mem_span_of_regionVirtualSlices,
        TNLean.PEPS.virtualBondSlice_mem_span_of_twoVertexRegionSlices}
    \leanok
    \uses{thm:peps_region_virtual_internal_bond,
        thm:peps_vertex_bond_coordinates,
        def:peps_virtual_bond_ground_space}
    Let $x$ be a global vector in independent virtual vertex coordinates.
    Suppose that, for every fixed complementary virtual configuration,
    its slice on $R$ belongs to the actual regional virtual bond range.
    Reorganize $x$ into its original edge-pair coordinates, obtaining $y$.
    For every internal edge $e$ of $R$ and every fixed configuration
    $\zeta$ of the other endpoint pairs, one has
    \begin{align*}
        y(\zeta[e\mapsto(a,b)])
            =c_\zeta\,\delta_{a,b}.
    \end{align*}
    Thus the corresponding one-bond slice belongs to the Bell line.
    In particular, the assertion applies to the two-vertex region
    formed by the endpoints of $e$. No bond dimension is assumed
    positive.
\end{theorem}

\begin{proof}\leanok
    Restrict a global configuration to $R$ and its complement.
    The regional support condition gives vanishing off the diagonal.
    Jointly permuting the two endpoint labels leaves the complementary
    configuration fixed because both endpoints belong to $R$.
    Transpositions therefore identify all diagonal coefficients.
    A coefficient function that vanishes off the diagonal and is
    constant on it is a scalar multiple of the Bell vector.
\end{proof}

\begin{theorem}[Regional virtual conditions imply all bond conditions]%
    \label{thm:peps_regional_virtual_conditions_bell}
    \lean{TNLean.PEPS.edgePair_mem_virtualBondGroundSpace_of_vertexVirtualParentGroundSpace}
    \leanok
    \uses{thm:peps_region_virtual_bell_slice,
        def:peps_virtual_bond_ground_space,
        def:peps_vertex_virtual_parent_ground_space}
    Let $(R_i)$ be a family of finite vertex regions such that the two
    endpoints of every edge lie together in some $R_i$.
    If every complementary slice of $x$ on every $R_i$ belongs to
    its actual regional virtual bond range, then its edge-pair
    coefficient vector belongs to the common Bell-bond ground space.
\end{theorem}

\begin{proof}\leanok
    For each edge choose a region containing both endpoints and apply
    the Bell-slice theorem. This gives every defining condition of
    the common Bell-bond space.
\end{proof}
